\clearpage\SourceLeaf{22}
\begin{center}
\bfseries ON THE THEORY OF CUBIC RESIDUES.
\end{center}



\begin{center}
[1.]
\end{center}

\[ p=mm+3nn \text{ a prime number.} \]

\noindent\hspace*{\parindent}For any integer indivisible by $p$, denoted by $x$, let
\[
\left[\dfrac{2mx}{p}\right]=\alpha,\qquad \left[\dfrac{(3n-m)x}{p}\right]=\beta,\qquad \left[\dfrac{(-3n-m)x}{p}\right]=\gamma.
\]
\noindent As is readily seen, this makes
\[
\dfrac{2mx}{p}-\alpha+\dfrac{(3n-m)x}{p}-\beta+\dfrac{(-3n-m)x}{p}-\gamma
\]
the sum of three proper fractions, so it must lie strictly between the limits
0 and 3. Since this same sum is evidently $=-\alpha-\beta-\gamma$, and hence an integer,
its only possible values are $+1$ and $+2$.
Thus either
\[ \alpha+\beta+\gamma=-1\quad\text{or}\quad \alpha+\beta+\gamma=-2. \]

\noindent It follows further that the three quantities $\alpha, \beta, \gamma$ cannot all be congruent
to one another modulo 3. Among the differences $\alpha-\beta$, $\beta-\gamma$,
$\gamma-\alpha$ there will therefore be at least one (and for that very reason at least one more,
